% =====================================================================
% Results — caching (co-headline). Phase 4 of RESEARCH_PLAN.md.
% =====================================================================
\section{Results: Per-Document Caching}
\label{sec:results-caching}

Per-document caching is the second half of the paper's contribution.
Under full causal attention, no document's KV block is reusable across
requests: $\mathrm{doc}_2$'s keys and values depend on $\mathrm{doc}_1$,
so changing any earlier document invalidates the prefix. Under a
block-diagonal mask with the restart position layout
(Section~\ref{sec:method-positions}), a document's KV block is a function
of that document alone. Blocks become keyable by document content hash,
order-independent, and reusable across requests.

A cache hit skips prefill entirely. This is the sense in which the
caching result is load-bearing: no kernel improvement to full-causal
prefill can match ``skip prefill'', and this speedup floor is in fact
the only wall-clock efficiency win this paper measures as working
end-to-end. The realization gap of Section~\ref{sec:results-efficiency}
is structural on the tested \texttt{sdpa} backend, so the caching path
carries the entire realized-speedup story rather than augmenting a
prefill-side win.

\subsection{Prerequisites}
\label{sec:results-caching-prereq}

Three conditions must hold before the throughput measurement is trusted,
and all three are settled positive in the current results.
\emph{Quality under restart} (Phase~4.1): on the composition probe with
the primary instruct model, \texttt{ddi\_qf + restart} sits within seed
noise of \texttt{ddi\_qf + sequential} on F1, presence, and loss
(headline: F1 $0.131$ vs $0.122$; presence $0.883$ vs $0.917$). Under
\texttt{full}, restart is not free --- it inflicts an F1 drop of
$\sim 58\%$ and a presence drop of $\sim 43$ points --- because full
attention exposes cross-document RoPE distances the model never saw at
training. The asymmetry is exactly the pairing the method was designed
to exploit: restart is quality-free precisely where DDI needs it to be.
\emph{Per-document cache correctness} (Phase~4.2): both deterministic
recompute and cache round-trip return $1.0$ across every doc$\times$record
tested. Per-document KV blocks are byte-stable at bf16.
\emph{Order invariance} (Phase~4.3): permuting the retrieved document
list under \texttt{ddi\_qf + restart} leaves logits invariant up to
floating-point precision --- the median max-delta collapses from
$\approx 0.31$ at bf16 to $\approx 0.017$ at fp32, matching the $28$-layer
accumulation floor and confirming that residual bf16 noise, not a real
position/mask leak, drove the earlier non-zero bf16 numbers. The
\texttt{sequential} arm's max-delta is unchanged by dtype ($\approx 3.6$)
because permuting doc order under sequential positions genuinely changes
token content and is expected to change logits. \todo{Bit-exact
invariance at \texttt{atol}$=10^{-3}$ is unattainable across $28$
decoder layers at any finite precision on current hardware; the
continuous max-delta comparison is the correct instrument for this
claim, and the strict binary threshold should be dropped from any table
that carries it.}

\subsection{Throughput under document reuse}
\label{sec:results-caching-throughput}

Figure~\ref{fig:f5} reports end-to-end throughput on a synthetic RAG
workload with document reuse (Zipf$(a{=}1.5)$ over a $32$-document
universe, $5$ documents per request, $200$ requests plus warmup, $3$
seeds). Three modes are compared: full-causal recompute (the status
quo), \texttt{ddi\_recompute} (block-diagonal masking with no cache),
and \texttt{ddi\_cache} (block-diagonal masking with a per-document KV
cache). The measured cache hit rate is $\approx 99\%$ at this reuse
distribution. \texttt{ddi\_cache} attains a median throughput of
$\approx 3.8$~QPS at a median request latency of $\approx 256$~ms,
against $\approx 2.2$~QPS at $\approx 459$~ms for full-causal and
$\approx 2.0$~QPS at $\approx 500$~ms for \texttt{ddi\_recompute} ---
a $\approx 1.77\times$ end-to-end throughput win over the status-quo
baseline. Variance is tight within each mode ($<1\%$ QPS spread across
seeds).

Ordering nuance worth naming: full-causal beats \texttt{ddi\_recompute}
by $\approx 9\%$ at these shapes because \texttt{ddi\_recompute} in the
current runner issues six separate small forward calls (five docs plus
query) while full-causal issues one concatenated forward, and per-call
Python and kernel-launch overhead outweighs the FLOP saving over
$\sim 500$-token doc blocks. \texttt{ddi\_recompute} is therefore not a
clean ``same compute, no cache'' control at these shapes: the caching
win is a caching win, not a per-call-compute win.
\todo{A batched multi-doc forward or a larger per-doc $L$ where FLOPs
dominate per-call overhead is needed to isolate the compute-vs-caching
axis cleanly; the current \texttt{ddi\_recompute} baseline should not
be read as pricing the mask's compute effect in isolation.}

\figureifexists{generated/F5.pdf}{%
    F5 --- throughput versus document reuse rate. One line per mode:
    full-causal recompute, DDI per-document recompute, DDI per-document
    cache.%
}{fig:f5}

The caching result survives failure modes elsewhere in the paper. If
the bridge-width sweep of Section~\ref{sec:results-quality-bridge}
shows that quality parity requires a wide bridge (small analytical
compute speedup), the caching result still holds --- the mask that
unlocks caching is the same document-block-diagonal mask, and the
bridge width does not affect cacheability of document blocks. Since
the realization gap of Section~\ref{sec:results-efficiency-gap} is
structural on the tested backend, the caching path is the only realized
wall-clock speedup this paper reports.
